% I. Описание на проблемната област
% Ориентировъчен обем: 6–7 стр.
\section{Описание на проблемната област}
\label{sec:domain}

\subsection{Изкуствени катерачни съоръжения и предаваните от тях товари}
\label{subsec:acs}

Изкуственото катерачно съоръжение (в текста нататък \term{съоръжение}, в изходния
текст и в интерфейса на системата с приетото съкращение ACS от англ.
\emph{artificial climbing structure}) е самостоятелна пространствена конструкция,
която не стои свободно, а се закрепва към носещата конструкция на сграда. Товарът,
който тя предава на сградата, има две различни по естество съставки:

\begin{itemize}[topsep=2pt,itemsep=1pt,leftmargin=*]
\item \term{постоянен товар} (означаван DL, от \emph{dead load}), собственото тегло на
      конструкцията, облицовката и хватките, което действа през целия живот на
      съоръжението и не се променя;
\item \term{полезен товар} (означаван LL, от \emph{live load}), товарът от катерача,
      включително динамичната съставка при падане и задържане на въжето, чиито
      стойности за съоръжения с осигурителни точки се определят по БДС EN
      12572-1~\cite{en12572_1}.
\end{itemize}

Двата вида съоръжения се разграничават по това дали имат осигурителни точки. При
\term{катерачната стена} катеренето е с въже през осигурителни точки, разположени на
едно или няколко нива по височина, и приложимата норма е EN
12572-1~\cite{en12572_1}. При \term{боулдер стената} осигурителни точки няма,
височината е малка и приземяването е върху постелка; приложимата норма е EN
12572-2~\cite{en12572_2}. Различието не е формално: то променя както големината на
полезния товар, така и точките, в които той се предава на сградата. За проекти извън
Европа системата използва американската спецификация на CWA~\cite{cwa2009}, което се
отразява и на частните коефициенти за товарите.

Разположението на товарите следва координатна система, чиито означения са възприети
без промяна от еталонните таблици и се използват както в изчислителното ядро, така и
в интерфейса на системата \tabref{tab:symbols}. Съществено за разбирането на цялата
работа е разграничението между двете групи величини. Величините с представка
\code{R} са реакциите върху \emph{една ос на съоръжението}: те описват какво се случва
в самото съоръжение и се използват при неговото оразмеряване. Величините с представка
\code{L} са товарите върху \emph{една колона на сградата} и вече отчитат разстоянието
между колоните. За отговора на въпроса „издържа ли сградата“ имат значение именно
величините \code{L}, което е причина интерфейсът да ги представя отделно и с по-силен
визуален акцент.

\begin{table}[htbp]
\centering
\caption{Означения на точките и величините в еталонните таблици}
\label{tab:symbols}
{\fontsize{10}{12}\selectfont
\begin{tabular}{@{}lL{9.6cm}@{}}
\toprule
\textbf{Означение} & \textbf{Значение} \\
\midrule
\code{RZ0} & вертикална съставка на реакцията в основата, на една ос на съоръжението \\
\code{RX0} & хоризонтална съставка на реакцията в основата, на една ос на съоръжението \\
\code{RX1}, \code{RX2}, \code{RX3} & хоризонтална реакция на ниво на закрепване 1, 2, 3 \\
\code{LX1}, \code{LX2}, \code{LX3} & хоризонтален товар върху колона на сградата на ниво 1, 2, 3, с отчетено разстояние между колоните \\
\code{DL} / \code{LL} & постоянен / полезен товар (наставка към всяко от горните) \\
$A$ & разстояние между колоните на сградата, m \\
$X$ & наклон на стената, разлика между горния и долния контур, m \\
$Z$ & височина на катерачната повърхност, m \\
\bottomrule
\end{tabular}}
\end{table}

Всички стойности в таблиците са \term{характеристични}, тоест ненамножени с частни
коефициенти. Преминаването към изчислителни стойности става чрез умножение с
коефициенти, които зависят от избраната нормативна система и са дадени в
\tabref{tab:factors}. Системата предлага това умножение като превключвател в
интерфейса, а не като отделен режим, за да остане видимо, че характеристичната
стойност е първичната.

\begin{table}[htbp]
\centering
\caption{Частни коефициенти за товарите по нормативна система}
\label{tab:factors}
{\fontsize{10}{12}\selectfont
\begin{tabular}{@{}lccL{5.2cm}@{}}
\toprule
\textbf{Система} & \textbf{DL} & \textbf{LL} & \textbf{Норма} \\
\midrule
Метрична (EU) & 1{,}35 & 1{,}50 & EN 12572-1 / EN 12572-2 \\
Имперска (US) & 1{,}20 & 1{,}60 & CWA 2009 \\
\bottomrule
\end{tabular}}
\end{table}

\subsection{Съдържание и структура на еталонните таблици}
\label{subsec:tables}

Изходният документ е работна книга на Microsoft Excel с макроси
(\code{Walltopia\_Preliminary\_Loads\_2026\_Rev.1.xlsm}) с обем 22{,}0 MB. От нея е
извлечен пълният набор от числа, който системата използва. Структурата, до която се
свежда съдържанието, е следната.

\begin{itemize}[topsep=2pt,itemsep=2pt,leftmargin=*]
\item \textbf{Катерачни стени.} 11 конфигурации, всяка определена от двойката
      (височина, брой нива на закрепване) и означавана с ключ от вида \code{W\_12\_3}.
      Височините са цели числа от 8 до 16 m, а броят нива е 1, 2 или 3, като не всяка
      комбинация е допустима: за височина 8 m съществуват схеми с 1 и с 2 нива, за
      9, 10 и 11 m само с 2 нива, за 12 m с 2 и с 3, а от 13 до 16 m само с 3 нива.
      Всяка конфигурация носи и височините на самите нива на закрепване, отделно в
      метрична и в имперска мярка.
\item \textbf{Боулдер стени.} 70 реда за две височини на катерачната повърхност,
      4 и 5 m, с една точка на закрепване.
\item \textbf{Свободни параметри.} За всяка конфигурация се задават разстояние между
      колоните $A \in \{4,5,6,7,8\}$ m и наклон $X \in \{1,\dots,7\}$ m. Общият брой
      редове за катерачни стени е 910.
\item \textbf{Стойности в един ред.} Всеки ред съдържа до дванадесет числа в метрична
      и същите дванадесет в имперска мярка, тоест еталонът пази двете мерни системи
      като независими набори, а не като преизчисление на едната от другата.
\end{itemize}

Последната подробност е проектно съществена и е причина системата да не прави
преобразуване на единици за силите. Ако имперските стойности се получаваха от
метричните чрез умножение, всяко закръгляне в еталона щеше да бъде възпроизведено
неточно. Затова изчислителното ядро \emph{избира} набор, а не преобразува
\tabref{tab:params}.

\begin{table}[htbp]
\centering
\caption{Обхват на еталонните данни}
\label{tab:params}
{\fontsize{10}{12}\selectfont
\begin{tabular}{@{}L{5.6cm}L{7.4cm}@{}}
\toprule
\textbf{Параметър} & \textbf{Стойности} \\
\midrule
Тип съоръжение & катерачна стена, боулдер стена \\
Мерна система & метрична (kN, m), имперска (lb, ft) \\
Височина, катерачна стена & 8, 9, 10, 11, 12, 13, 14, 15, 16 m \\
Височина, боулдер стена & 4, 5 m \\
Брой нива на закрепване & 1, 2, 3 (според височината) \\
Разстояние между колоните $A$ & 4, 5, 6, 7, 8 m \\
Наклон $X$ & 1, 2, 3, 4, 5, 6, 7 m \\
Конфигурации на катерачни стени & 11 \\
Редове, катерачни стени & 910 \\
Редове, боулдер стени & 70 \\
\bottomrule
\end{tabular}}
\end{table}

\subsection{Недостатъци на съществуващото решение}
\label{subsec:asis}

Днес описаните таблици се използват във вида, в който са произведени, тоест като
работна книга, разпращана по електронна поща. Този начин на работа поражда пет
различни по естество проблема \tabref{tab:asis}. Те не са недостатъци на самите
таблици: съдържанието им е вярно и не се оспорва. Всички произтичат от носителя и
от начина на разпространение, което е и причината задачата да бъде решима с програмна
система, без да се пипа конструктивното изчисление.

\begin{table}[htbp]
\centering
\caption{Проблеми при работата с еталонната работна книга}
\label{tab:asis}
{\fontsize{10}{12}\selectfont
\begin{tabular}{@{}L{3.3cm}L{5.0cm}L{4.6cm}@{}}
\toprule
\textbf{Проблем} & \textbf{Проявление} & \textbf{Изискване, което го покрива} \\
\midrule
Зависимост от среда & изисква настолен Microsoft Excel с разрешени макроси & NFR7, достъп през браузър \\
Неконтролирани копия & файлът се разпраща и се съхранява в неизвестен брой версии & един източник, обслужван от услугата \\
Липса на следа & не се знае кой, кога и с какви входни данни е получил дадено число & FR8, запазване на проекти \\
Неизползваемост извън офиса & 22 MB работна книга не се отваря удобно от телефон & NFR4, FR11 \\
Риск от грешно четене & стойността е на пресечната точка на три избора; грешният избор не се сигнализира & FR1, ограничаване на избора до допустимото \\
\bottomrule
\end{tabular}}
\end{table}

Последният ред от \tabref{tab:asis} е и най-съществен от инженерна гледна точка.
Работната книга е пасивна: тя показва всички редове и оставя на четящия да избере
верния. Ако избере схема с 3 нива за стена от 9 m, каквато таблиците не описват, той
или ще намери празна клетка, или, при невнимание, ще прочете съседна стойност.
Програмната система може да не допусне този избор изобщо, което превръща цял клас
човешки грешки в невъзможно състояние. Това е основната функционална разлика между
двете и е причината §\ref{subsec:core} да е посветен именно на нея.

\subsection{Потребители и сценарии на употреба}
\label{subsec:users}

Системата обслужва две различни по цел групи, чиито изисквания не съвпадат.

\textbf{Външен потребител} (проектант, архитект или инвеститор, който обмисля
съоръжение в конкретна сграда) търси бърз отговор на въпроса дали идеята е
осъществима. За него са съществени малкият брой стъпки до резултата, разбираемото
означение на величините и възможността да провери получения товар срещу носимостта
на съществуващата конструкция. Той работи от произволно устройство, често не от
работно място с настолен компютър.

\textbf{Вътрешен потребител} (търговският и инженерният отдел на Walltopia) има нужда
от следа: кой клиент какво е сметнал, с какви входни данни, кога, и какъв въпрос е
задал по този конкретен случай. Днес тази следа не съществува, защото изчислението се
прави в копие на работната книга върху чужда машина.

От двете групи произтичат трите основни сценария, около които е изградена системата:

\begin{enumerate}[topsep=2pt,itemsep=2pt,leftmargin=*]
\item \textbf{Еднократна проверка.} Потребителят избира мерна система, тип, височина,
      схема на закрепване, разстояние между колоните и наклон, и получава таблицата с
      товарите заедно със схема на съоръжението. Не е необходима регистрация.
\item \textbf{Запазване като проект.} След регистрация изчислението се запазва с име,
      етикети и произволен брой двойки „свойство: стойност“, и може да бъде отворено и
      променено по-късно. Списъкът с проекти се търси и филтрира.
\item \textbf{Запитване към инженерния отдел.} От запазен проект се изпраща въпрос,
      към който автоматично се прикачват входните данни и обобщението на резултата, за
      да не се налага потребителят да ги преписва.
\end{enumerate}

\subsection{Функционални изисквания}
\label{subsec:fr}

От сценариите са изведени функционалните изисквания в \tabref{tab:fr}. Всяко от тях е
проследимо до реализация в Раздел~\ref{sec:implementation}, до проверка в
§\ref{subsec:method} и до описание за потребителя в Раздел~\ref{sec:userguide}.

\begin{table}[htbp]
\centering
\caption{Функционални изисквания}
\label{tab:fr}
{\fontsize{10}{12}\selectfont
\begin{tabular}{@{}lL{11.4cm}@{}}
\toprule
\textbf{№} & \textbf{Изискване} \\
\midrule
FR1 & Избор на мерна система, тип съоръжение, височина, схема на закрепване, разстояние между колоните и наклон, с автоматично ограничаване до допустимите стойности при всяка промяна. \\
FR2 & Извеждане на характеристичните товари за избраната комбинация, разделени на реакции върху ос на съоръжението и товари върху колона на сградата. \\
FR3 & Превключване към изчислителни стойности чрез частните коефициенти на избраната нормативна система. \\
FR4 & Проверка на получения товар срещу въведена от потребителя носимост, с еднозначна оценка „приложимо“ или „надвишена носимост“. \\
FR5 & Извеждане на две алтернативни схеми на закрепване, пряко закрепване в отделни точки и закрепване през греди между колоните на сградата. \\
FR6 & Мащабируема схема на съоръжението, която се преизчислява от текущия избор. \\
FR7 & Регистрация, вход и изход на потребител. \\
FR8 & Запазване, отваряне, редактиране и изтриване на проект с име, етикети и потребителски свойства. \\
FR9 & Търсене, филтриране по етикет и подреждане на списъка с проекти. \\
FR10 & Изпращане на запитване към инженерния отдел, обвързано с конкретен проект. \\
FR11 & Достъп до приложното ръководство и до типовите детайли за закрепване. \\
\bottomrule
\end{tabular}}
\end{table}

\subsection{Нефункционални изисквания}
\label{subsec:nfr}

Нефункционалните изисквания са дадени в \tabref{tab:nfr}. Две от тях заслужават
отделен коментар, защото определиха проектни решения, а не само качество на
изпълнението.

NFR2 изисква калкулаторът да работи и когато базата от данни е недостъпна. Причината
е, че товарните таблици не са потребителски данни: те са едни и същи за всички и не
зависят от вход в системата. Изискването е реализирано като разделяне на маршрутите,
описано в §\ref{subsec:server}.

NFR4 изисква интерфейсът да остава използваем при всяка ширина на изгледа от 320 до
2560 px без хоризонтално превъртане на страницата. Изискването е формулирано като
механично проверимо свойство именно за да може да бъде проверено автоматично; как
точно, е описано в §\ref{subsec:howcheck}.

\begin{table}[htbp]
\centering
\caption{Нефункционални изисквания}
\label{tab:nfr}
{\fontsize{10}{12}\selectfont
\begin{tabular}{@{}lL{11.4cm}@{}}
\toprule
\textbf{№} & \textbf{Изискване} \\
\midrule
NFR1 & Изчислителният резултат съвпада с еталонната работна книга за всяка допустима входна комбинация. \\
NFR2 & Калкулаторът остава работоспособен при недостъпна база от данни; недостъпни са само функциите, свързани с потребителски профил. \\
NFR3 & Интерактивните точки в схемата на съоръжението и в конфигуратора на закрепването са достъпни от клавиатура (фокусируеми, с роля на бутон) и носят описателен \code{aria-label}. \\
NFR4 & Липса на хоризонтално препълване на страницата при ширина на изгледа от 320 до 2560 px. \\
NFR5 & Паролите се съхраняват само като хеш; сесията на уеб клиента не е достъпна от скрипт на страницата. \\
NFR6 & Ограничаване на честотата на заявките към входните точки за удостоверяване. \\
NFR7 & Уеб клиентът се доставя без стъпка на компилация, за да може да бъде публикуван като статични файлове. \\
NFR8 & Всяка страница с резултати съдържа видимо указание, че товарите са предварителни. \\
\bottomrule
\end{tabular}}
\end{table}

\subsection{Ограничения на предметната област}
\label{subsec:limits}

Системата дава \term{предварителни} товари. Това не е предпазна формулировка, а точно
описание на обхвата: таблиците са изготвени за типови конфигурации при типови
допускания и служат за преценка на осъществимост, не за изпълнителен проект. Три
неща остават извън тях и не могат да бъдат получени от системата:

\begin{enumerate}[topsep=2pt,itemsep=1pt,leftmargin=*]
\item оразмеряването на самите закрепвания, което зависи от материала и състоянието
      на съществуващата конструкция;
\item всяко въздействие извън собственото тегло и товара от катерача, например
      сеизмично или температурно;
\item конфигурациите извън обхвата на \tabref{tab:params}, за които системата
      съзнателно не екстраполира, а съобщава, че за тази комбинация няма ред в
      таблицата.
\end{enumerate}

Последното е проектно решение с последица за интерфейса: липсата на данни се
съобщава като липса, а не се запълва с изчислена по съседни редове стойност.
Причината е, че мълчаливата екстраполация би произвела число, което изглежда като
резултат от таблица, но не е.
